% Generated by experiments/render_results.py; do not edit numerical tables by hand.

\subsection{Study 1: fault detection and false alarms}

Table~\ref{tab:faults} reports alarms for each corruption. The six explicit output violations were detected in every evaluation case. Both well-formed ensemble mutations were rejected by the covariance check in every case.

\begin{table}[htbp]\centering\scriptsize

\caption{Alarm counts and Wilson 95\% intervals (percent). Covariance is not run on malformed outputs; --- denotes an undefined result, not a negative result.}

\label{tab:faults}

\begin{tabular}{@{}lccc@{}}\toprule

Case & Point-MSE alarm & Output contract & Covariance \\ \midrule

Clean control & 5/64 [3.4, 17.0] & 0/64 [0.0, 5.7] & 6/64 [4.4, 19.0] \\

Index shift & 5/64 [3.4, 17.0] & 64/64 [94.3, 100.0] & --- \\

Target-label swap & 5/64 [3.4, 17.0] & 64/64 [94.3, 100.0] & --- \\

Horizon truncation & 64/64 [94.3, 100.0] & 64/64 [94.3, 100.0] & --- \\

Nonfinite point & 64/64 [94.3, 100.0] & 64/64 [94.3, 100.0] & --- \\

Quantile crossing & 5/64 [3.4, 17.0] & 64/64 [94.3, 100.0] & --- \\

Sample-axis transpose & 5/64 [3.4, 17.0] & 64/64 [94.3, 100.0] & --- \\

Variance shrinkage & 5/64 [3.4, 17.0] & 0/64 [0.0, 5.7] & 64/64 [94.3, 100.0] \\

Dependence shuffle & 5/64 [3.4, 17.0] & 0/64 [0.0, 5.7] & 64/64 [94.3, 100.0] \\

\bottomrule\end{tabular}\end{table}

The clean covariance check rejected 6/64 cases (9.4\%; Wilson interval [4.4, 19.0]\%). The observed rate is above the nominal 5\% per-check level, with an interval that includes 5\%. No threshold was changed to remove these false alarms. This run does not establish precise calibration or a family-wise error rate.

The point-MSE cutoff was 4.1853, the empirical 95th percentile from 16 development cases. It alarmed on 5/64 clean cases. It detects invalid point arrays, but ignores index labels, target names, quantiles, and trajectory samples. Where a mutation leaves the point values unchanged, its alarms are identical to the clean alarms and supply no mutation-specific evidence.

Disabling an output mutation restored the exact clean output, including its sample array. Repaired outputs had zero contract violations and the same 6/64 covariance alarms as the clean control. These paired rechecks confirm removal of the corruption; they do not imply that every repaired case receives a statistical pass.

\subsection{Study 1: marginal invariance and joint-score response}

Table~\ref{tab:dependence} varies the proportion of ensemble rows selected for coordinate-wise permutation. Every empirical marginal is preserved at every strength, as checked directly in the runner. The maximum absolute CRPS difference across all cases and strengths was $8.88\times10^{-16}$. Energy and variogram scores both worsened on average at full permutation; their paired intervals exclude zero for this setting. Their raw effect magnitudes have different units and must not be used to rank their sensitivity.

\begin{table}[htbp]\centering\small

\caption{Mutated minus clean score, mean and paired 95\% t interval over 64 independent cases. CRPS differences are floating-point roundoff. Intervals are exploratory and not multiplicity-adjusted.}

\label{tab:dependence}

\begin{tabular}{@{}rccc@{}}\toprule

Strength & CRPS mean & Energy & Variogram \\ \midrule

0.00 & $0$ & +0.0000 [+0.0000, +0.0000] & +0.0000 [+0.0000, +0.0000] \\

0.25 & $-7.81\times10^{-18}$ & +0.0150 [+0.0037, +0.0263] & +0.0069 [+0.0019, +0.0120] \\

0.50 & $6.07\times10^{-18}$ & +0.0389 [+0.0209, +0.0570] & +0.0194 [+0.0102, +0.0286] \\

0.75 & $-8.67\times10^{-19}$ & +0.0679 [+0.0456, +0.0902] & +0.0295 [+0.0181, +0.0408] \\

1.00 & $-4.34\times10^{-18}$ & +0.0767 [+0.0534, +0.1001] & +0.0351 [+0.0230, +0.0472] \\

\bottomrule\end{tabular}\end{table}

\subsection{Study 1: capability profiles}

The univariate capability sweep uses 16 evaluation seeds, three innovation scales, five kinds of series, and three unmodified models. Table~\ref{tab:capability} scores each model against the answer key at the lowest and highest scale: 100 times the unavoidable error over the sum of the unavoidable and avoidable errors (Section~\ref{sec:references}). The full three-scale sweep is in \code{figures/capability-curves.pdf} and \code{capability-summary.csv} of the reproducibility artifact. The fitted AR(1) nails white noise and mean reversion and fails on the trend at both scales. For the zero-drift random walk, the last value is the answer key. Noise matters where a fixed-size pattern competes with it: on the seasonal cycle the fitted AR(1) scores 9 at low noise and 91 at high noise, where the cycle is mostly buried and ignoring it costs little.

\begin{table}[htbp]\centering\small

\caption{Score out of 100 (100 matches the answer key) at low ($\sigma=0.25$) and high ($\sigma=2$) innovation noise, over 16 cases each. N nails it, G good, S struggles, F fails. White noise, mean reversion and the random walk score the same at both levels: with a zero mean and no drift, a larger $\sigma$ only rescales the series. The package's scorecard adds measurement noise to tell those levels apart.}

\label{tab:capability}

\begin{tabular}{@{}lrrrrrr@{}}\toprule
 & \multicolumn{2}{c}{Last value} & \multicolumn{2}{c}{Seasonal naive (12)} & \multicolumn{2}{c}{Fitted AR(1)} \\
\cmidrule(lr){2-3}\cmidrule(lr){4-5}\cmidrule(lr){6-7}
Kind of series & Low & High & Low & High & Low & High \\ \midrule

Noise around a fixed level & 51 S & 51 S & 55 S & 55 S & 99 N & 99 N \\

Straight-line trend & 17 F & 50 S & 4 F & 47 S & 39 F & 18 F \\

Repeating seasonal cycle & 5 F & 46 S & 55 S & 55 S & 9 F & 91 N \\

Mean reversion (AR(1)) & 68 S & 68 S & 40 S & 40 S & 95 N & 95 N \\

Random walk & 100 N & 100 N & 37 F & 37 F & 87 G & 87 G \\

\bottomrule\end{tabular}\end{table}

\subsection{Study 1: an inverse-scaling repair}

A separate predictor standardizes its training context and forecasts the mean in standardized units. The fault omits the inverse transform; the repair applies the training mean and scale. The historical-mean predictor is shift equivariant, so this particular metamorphic relation is a valid declared contract. On white-noise cases with mean 5 and innovation scale 2, the relation failed for 64/64 faulty fits (Wilson 95\% interval [94.3, 100.0]\%) and 0/64 repaired fits ([0.0, 5.7]\%). The paired repaired-minus-faulty MAE difference was -3.4711 [-3.6333, -3.3090] (95\% t interval).

We then apply the same known omission and repair to AirPassengers, the monthly 1949--1960 international-airline passenger series in thousands~\cite{rAirPassengers}. All scaling statistics are fitted on each training prefix. The four ordered origins each forecast the next 12 months; their results are reported separately because they are not independent series replicates.

\begin{table}[htbp]\centering\small

\caption{AirPassengers inverse-scaling illustration. MAE is in thousands of passengers. The declared shift relation fails before and passes after repair at all four origins.}

\label{tab:real}

\begin{tabular}{@{}lrrr@{}}\toprule

Last training month & Faulty MAE & Repaired MAE & Difference \\ \midrule

1956-12 & 368.42 & 154.71 & -213.71 \\

1957-12 & 381.00 & 150.10 & -230.90 \\

1958-12 & 428.33 & 182.43 & -245.91 \\

1959-12 & 476.17 & 213.67 & -262.49 \\

\bottomrule\end{tabular}\end{table}

Repair improves MAE at every displayed origin, but the repaired historical mean remains a poor model of this growing seasonal series.
